\documentclass[11pt]{article}
\usepackage{fontspec}
\setmainfont{Times New Roman}
\setsansfont{Arial}
\setmonofont{Consolas}
\usepackage{amsmath,amssymb}
\usepackage{geometry}
\usepackage{microtype}
\geometry{a4paper,margin=3cm}
\setlength{\parindent}{1.25em}
\setlength{\parskip}{0pt}
\emergencystretch=8em
\allowdisplaybreaks
\newcommand{\pair}[2]{\left(#1\middle|#2\right)}
\newcommand{\eps}{\varepsilon}

\begin{document}

\section*{\S\ 5. Тотожності розкладання в явній формі.}

Якщо у випадку 4) в (25.) спеціалізувати $\tau$ до $\sigma+\rho$, то маємо:
\begin{equation}
\begin{aligned}
 \pair{q_\rho A^\sigma}{y^{(1)}\cdots y^{(\rho+\sigma)}}
 &=\sum_{1\le i_1<\cdots<i_\rho\le \rho+\sigma}
 \eps_{i_\rho}\pair{q_\rho}{y^{(i_1)}\cdots y^{(i_\rho)}}\\
 &\qquad\cdot
 \pair{A^\sigma}{y^{(i_{\rho+1})}\cdots y^{(i_{\rho+\sigma})}},
\end{aligned}
\tag{32}
\end{equation}
де $\eps_{i_\rho}=(-1)^{\delta_{i_\rho}}$, $\delta_{i_\rho}=\rho(\rho+1)/2+i_1+\cdots+i_\rho$. Це загальніша форма теореми про добуток для матриць, ніж (16.) і (17.); вона випливає з лапласового розкладу детермінанта продуктової матриці
\[
 \sum\pm(v^{(1)}y^{(1)})\cdots(v^{(\rho)}y^{(\rho)})
 (a^{(1)}y^{(\rho+1)})\cdots(a^{(\sigma)}y^{(\rho+\sigma)}).
\]
Отже, загальніша теорема про добуток складена з тотожностей (20.), відповідно (21.), і тим самим пояснено виокремлення спеціальної форми (16.), (17.).

Співвідношення (32.) тепер дає засіб для явного подання тотожностей розкладання. Для цього розглянемо форми, що входять у (31.), як основні форми; тобто виконаємо підстановку (11.), яка нічого не змінює в явному поданні через рядки $A,B$, бо за (8.) нововведені рядки $R$ виражаються через самі рядки $A,B$, а не через їхні часткові рядки $a^{(1)}a^{(2)}\ldots$, $b^{(1)}b^{(2)}\ldots$. Нехай, отже,
\begin{equation}
 \pair{q_{\rho-\alpha}B^{n-\rho-\sigma+\tau}}{A_{n-\sigma}p_{\tau-\alpha}}
 =\pair{R_{\rho-\alpha}p_{n-\rho+\alpha}}{R^{\tau-\alpha}p_{\tau-\alpha}},
\tag{33}
\end{equation}
тоді для окремої суми з (25.), що відповідає індексу $\alpha$, дістаємо:
\begin{equation}
\begin{aligned}
&\sum \eps_{i_\alpha}
\pair{R_{\rho-\alpha}p_{n-\rho}y^{(i_1)}\cdots y^{(i_\alpha)}}
{R^{\tau-\alpha}y^{(i_{\alpha+1})}\cdots y^{(i_\tau)}}\\
&\quad=\sum \eps_{i_\alpha}
\pair{[R_{\rho-\alpha}p_{n-\rho}]^\alpha
y^{(i_1)}\cdots y^{(i_\alpha)}}
{R^{\tau-\alpha}y^{(i_{\alpha+1})}\cdots y^{(i_\tau)}}\\
&\quad=\eps\pair{[R_{\rho-\alpha}p_{n-\rho}]^\alpha R^{\tau-\alpha}}{p_\tau}
 =\eps\pair{R^{\tau-\alpha}q_{n-\tau}}{R_{\rho-\alpha}p_{n-\rho}}.
\end{aligned}
\tag{34}
\end{equation}
Цим справді задано явний підсумковий вираз, а симетрична форма, у якій входять $q_\rho$ і $p_\tau$, показує, що (30.) приводить до того самого підсумкового виразу, який, отже, не залежить від початкового розкладання. Різниця між (25.) і (30.) існує лише доти, доки самі рядки $y$ і $v$ мають самостійне значення, тобто доки тотожність розкладання за нашим означенням переходить у розкладну тотожність. Щоб перейти до тотожностей з більшою кількістю символьних рядків, треба за \S\ 3 (кінець) лише розкласти рядок $q_\rho$ у $q_{\rho_1}C^{\sigma_1}$, відповідно $p_\tau$ у $p_{\tau_1}C_{\tau-\tau_1}$. Тоді виникають вирази:
\begin{equation}
 \pair{q_{\rho_1}C^{\sigma_1}}
 {[R^{\tau-\alpha}q_{n-\tau}]_\lambda R_{\rho+\sigma_1-\alpha-\lambda}},
 \qquad
 \pair{[R_{\rho-\alpha}p_{n-\rho}]^\lambda R^{\tau-\alpha}}
 {p_{\tau_1}C_{\tau-\tau_1}},
\tag{35}
\end{equation}
а вони мають точно той самий тип, що й у випадку двох символьних рядків, і тому за підстановки, відповідної до (33.), знову допускають явне подання; повторенням дістаємо явне подання для довільної кількості рядків. Варто ще зауважити, що перша і остання суми в (25.), відповідно (30.), дають явне подання без застосування (33.), лише за формулою (32.), або взагалі зводяться до одного члена, який явно містить усі рядки. Співвідношення, що виникає з (32.) через розкладання $q_\rho$ у $q_{\rho_1}C^{\sigma_1}$ тощо, можна розглядати як теорему про добуток у її найзагальнішій формі.

Підсумовуючи, покажемо:

\emph{Теорема III:} Тотожностями
\begin{equation}
 \pair{q_\rho A^\sigma}{p_\tau B_{\rho+\sigma-\tau}}
 =\sum_{\alpha=\max(0,\tau-\sigma)}^{\min(\tau,\rho)}
 \eps\pair{R^{\tau-\alpha}q_{n-\tau}}{R_{\rho-\alpha}p_{n-\rho}},
\tag{36}
\end{equation}
де рядки $R$ визначаються через (33.), і тотожностями, що виникають з них через розкладання $q_\rho,p_\tau$ у $q_{\rho_1}C^{\sigma_1}$ тощо, вичерпується вся сукупність нерозкладних тотожностей.

Справді, тотожності, що виникають так, є найзагальнішими у своєму роді, бо окремі рядки вже не підлягають жодним обмеженням щодо ваги й кількості, як це ще було у випадку (20.), (21.). Але між цими рядками не існує й жодних інших тотожностей: адже при розкладанні рядка $p_\tau$ у $y^{(1)}\cdots y^{(\tau)}$ найзагальніші тотожності складені з (20.), отже, за \S\ 3 (кінець), з (20a.); причому композиція, обговорена після (16.), є саме тією, яку проведено в \S\ 4. Перехід до довільної кількості рядків, як показує обговорення (20a.), дає застосування тієї самої операції до виразів, що виникають через розкладання $q_\rho$ у $q_{\rho_1}C^{\sigma_1}$ тощо. Звідси також випливає, що прямий перехід від (20.) замість (20a.) не дає нічого нового; це легко перевірити обчисленням, якщо один раз спеціалізувати рядок $q_\rho$ у (25.), а другий раз виконати перехід від (20.) за методом \S\ 4. В обох випадках виникають ті самі суми, які через застосування загальної теореми про добуток (див. вище) переходять у тотожності, що випливають із (36.). Тотожності (20.), (21.) відповідають спеціальним випадкам $\tau=1$, $\rho=1$; тоді з'являються лише по дві окремі суми, а отже, відповідно до наведеного вище зауваження, явне подання без підстановки (33.), що узгоджується з попереднім.

На завершення коротко згадаємо ще два спеціальні випадки тотожності розкладання. Якщо в (31.) покласти $\tau=\sigma$, $\rho=n-\sigma$ і позначити рядок $p_\tau$ як $p'_\sigma\sim q'_{n-\sigma}$, то маємо:
\begin{equation}
\begin{aligned}
\pair{A^\sigma}{p_\sigma}\pair{B^\sigma}{p'_\sigma}
-\pair{A^\sigma}{p'_\sigma}\pair{B^\sigma}{p_\sigma}
=\sum_{\alpha=1}^{\min(\sigma,n-\sigma)}
\Pi\pair{q_{n-\sigma-\alpha}B^\sigma}{A_{n-\sigma}p_{\sigma-\alpha}}.
\end{aligned}
\tag{37}
\end{equation}
Це формула Клебша\footnote{Clebsch, a. a. O., S. 25--32. Клебш доводить, що окремі члени правих частин залежать лише від рядків $A,B$; явне подання отримується застосуванням (32.) до його виразів $\Phi_h$.} для розвинення форми з двома когредієнтними рядками $p_\sigma,p'_\sigma$ за елементарними коваріантами. Елементарні коваріанти, що належать до форми $\pair{A^\sigma}{p_\sigma}^m\pair{B^\sigma}{p'_\sigma}^n$, тоді такі:
\begin{equation}
 \left\{\sum_{\alpha=1}^{\min(\sigma,n-\sigma)}
 \eps\pair{q_{n-\sigma-\alpha}B^\sigma}{A_{n-\sigma}p_{\sigma-\alpha}}
 \right\}^{\rho}
 \pair{A^\sigma}{p_\sigma}^{m-\rho}\pair{B^\sigma}{p'_\sigma}^{n-\rho},
 \qquad (\rho=0,1,\ldots,m,n).
\tag{38}
\end{equation}
Вони, як ми коротко казатимемо, виникають з основної форми через ``когредієнтне згортання з дефектом $\rho$'' (пор. \S\ 2).

По-друге, якщо покласти $\tau=\sigma-\lambda$, $\rho=n-\sigma-\lambda$; $B^\sigma=A^\sigma$, то маємо:
\begin{equation}
\begin{aligned}
\pair{q_{n-\sigma-\lambda}A^\sigma}{A_{n-\sigma}p_{\sigma-\lambda}}
&=(-1)^\lambda
\pair{q_{n-\sigma-\lambda}A^\sigma}{A_{n-\sigma}p_{\sigma-\lambda}}\\
&\quad+(-1)^{(n-\sigma)(\sigma-\lambda)}
 \sum_{\alpha=1}^{\min(\sigma-\lambda,n-\sigma-\lambda)}
 \Pi\pair{q_{n-\sigma-\lambda-\alpha}A^\sigma}
 {A_{n-\sigma}p_{\sigma-\lambda}}.
\end{aligned}
\tag{39}
\end{equation}
Отже, як зазначено в \S\ 2, умови (12.), що характеризують символьні рядки, для непарного дефекту $\lambda$ є наслідком умов вищого дефекту. Для лінійної форми $\pair{A^\sigma}{p_\sigma}=\pair{B^\sigma}{p_\sigma}$ з (39.) випливає, що її нередуковні інваріантні утворення, квадратичні за коефіцієнтами, зводяться до утворень парного дефекту; це узгоджується з відомим фактом, що лінійна форма в бінарній і тернарній областях не має інваріантних утворень.

Зауважимо, що загальні тотожності спочатку були виведені за припущення, що окремі рядки задовольняють умови (12.). Але оскільки ці окремі рядки входять у тотожності лише лінійно, останні справджуються також для загальніших рядків, які не задовольняють умови (12.).

\end{document}
